\section{Alina Kopasava}
\label{sec:wklad-kopasava}

Zakres prac obejmował całą aplikację kliencką: jej architekturę, wszystkie widoki oraz scenę trójwymiarową.

Zaprojektowany został podział kodu na cztery warstwy opisany w podrozdziale~\ref{sec:implementacja-frontendu}, oddzielający typy dziedzinowe od logiki aplikacji, integracji sieciowych i prezentacji. W tym układzie powstały wszystkie strony aplikacji: strona główna z formularzem kontaktowym, katalog materiałów, formularze rejestracji, logowania i odzyskiwania hasła, panel administratora z wykazem użytkowników, kart zamówień i wiadomości oraz uproszczony widok montażowy przystosowany do ekranu telefonu. Dostępem do widoków zastrzeżonych steruje komponent sprawdzający rolę użytkownika przed wyświetleniem strony.

Najobszerniejszą częścią prac był konfigurator produktu wraz ze sceną trójwymiarową. Do tego zakresu należą generowanie geometrii trzynastu kształtów tablicy, proceduralne tworzenie tekstur kamienia, nakładanie inskrypcji z zewnętrznych krojów pisma, kadrowanie wgranej fotografii portretowej do ramki na tablicy oraz bieżące przeliczanie ceny według wzoru~(\ref{eq:cena}). Zrealizowana została również wielojęzyczność interfejsu w trzech wersjach oraz prezentacja cen w walucie zależnej od wybranego języka, wraz z obsługą sytuacji, w której kurs jest niedostępny.

Po stronie integracji z serwerem powstał moduł skupiający cały ruch sieciowy aplikacji klienckiej, przekształcający odpowiedzi błędne w wyjątki i rozpoznający przekroczenie limitu liczby żądań. W obszarze testów przygotowane zostały testy jednostkowe funkcji dziedzinowych aplikacji klienckiej oraz testy komponentowe stron wykonywane w środowisku jsdom. Do zakresu prac należało ponadto opracowanie dokumentacji projektowej.
